\chapter{Acknowledgements}


From topic selection, proposal, and mid-term progress to final completion, this graduation project went through repeated requirement adjustments, architecture iterations, deployment troubleshooting, and thesis revision. The successful completion of this work would not have been possible without the continuous guidance of my supervisor in both overall direction and detailed improvements. At each key stage, my supervisor provided specific and practical suggestions, helping me transform initial ideas into concrete implementation goals and deepening my understanding of usability and interpretability in real teaching scenarios.

At the same time, course training and laboratory projects in the school provided a strong foundation for this project, enabling me to integrate database design, backend development, API integration, and deployment operations into a fully runnable system. During development and testing, I also benefited from classmates' feedback on usage experience, function behavior, and issue reproduction. Their practical input helped identify and locate many problems that were difficult to discover through self-testing alone.

For me personally, this graduation project is not only an academic requirement, but also a complete engineering practice experience. I learned how to balance feature implementation and system stability under limited time, how to locate problems in a layered manner when facing continuous errors, and how to translate technical implementation into solutions acceptable for real teaching contexts. Finally, I would like to express my sincere gratitude to all teachers and classmates who supported this project.

